\documentclass[aps,prl,reprint,amsmath,amssymb,nofootinbib,raggedbottom]{revtex4-2}
\usepackage[T1]{fontenc}
\usepackage{tnnotes}
\usepackage[colorlinks=true,citecolor=blue,linkcolor=blue,urlcolor=blue]{hyperref}
\hypersetup{pdftitle={Scope of the Two-Purity Method}}

\newtheorem{corollary}{Corollary}

% ------------------------------------------------ diagrams specific to this note
% One bond event: C on the left site, Theta(C) on its mirror image; the dashed
% line is the mirror, the double lines are the two sites' imaginary-time legs.
\newcommand{\bondform}[2]{\tikz[tn,tn centered]{
 \draw[phys] (0,-1)--(0,1); \draw[phys] (3.2,-1)--(3.2,1);
 \draw (0,0)--(3.2,0);
 \draw[black!45,densely dashed,line width=.35pt] (1.6,-1.15)--(1.6,1.15);
 \node[xop,fill=white] at (0,0) {$#1$};
 \node[xop,fill=white] at (3.2,0) {$#2$};}}
% A transfer vertex with named virtual legs.
\newcommand{\Xab}[2]{\tikz[tn]{\draw (-.95,0)--(.95,0); \node[xvertex] at (0,0) {};
 \node[anchor=east] at (-1,0) {$#1$}; \node[anchor=west] at (1,0) {$#2$};}}
% One site of the transfer vertex, opened in imaginary time: the physical
% trace (closed double loop) with the left-bond insertions #1 and the
% right-bond insertions #2.
\newcommand{\eventsite}[2]{\tikz[tn,tn centered]{
 \begin{pgfonlayer}{background}
 \draw[phys,ring] (0,-1.9)--(0,1.9)--(1.9,1.9)--(1.9,-1.9)--cycle;
 \draw (-1.9,1)--(0,1); \draw (-1.9,-1)--(0,-1);
 \draw[preaction={draw,white,line width=2.5pt}] (0,0)--(2.6,0);
 \end{pgfonlayer}
 \node[xop,fill=white] at (0,1) {$#1$};
 \node[xop,fill=white] at (0,-1) {$#1$};
 \node[xop,fill=white] at (0,0) {$#2$};
 \node[anchor=east] at (-1.9,0) {$a$}; \node[anchor=west] at (2.6,0) {$b$};}}
% The reflection-positivity inequality: a spacetime torus cut by the mirror,
% A on the left half, theta(A) on its image.
\newcommand{\rpcut}{\tikz[tn,tn centered]{
 \fill[mpoblue!12] (0,0) rectangle (4.4,1.8);
 \draw[black!60,line width=.35pt] (0,0) rectangle (4.4,1.8);
 \draw[black!45,densely dashed,line width=.35pt] (2.2,-.25)--(2.2,2.05);
 \node[xop,fill=white] at (1.1,.9) {$A$};
 \node[xop,fill=white] at (3.3,.9) {$\theta(A)$};}}
% delta(eps) = [(1-eps)^-3 - 1]^2/2 against the diagonal, 0 <= eps <= .16;
% they cross at eps* = .127.
\newcommand{\deltacurve}{\tikz[baseline=0pt,x=6.5em,y=6.5em,line width=.45pt]{
 \draw[black!35,line width=.3pt] (0,0)--(.16,0) (0,0)--(0,.16);
 \draw[black!45,densely dotted] (0,0)--(.16,.16);
 \draw[mpored,domain=0:.1405,samples=30,smooth]
   plot (\x,{.5*(1/((1-\x)*(1-\x)*(1-\x))-1)*(1/((1-\x)*(1-\x)*(1-\x))-1)});
 \fill (.1273,.1273) circle[radius=.035em];}}

\begin{document}
\title{Scope of the Two-Purity Method}
\date{October 8, 2026}
\begin{abstract}
The two-purity method establishes a uniform spectral gap from bond reflection positivity and certified finite-scale purity bounds. Bond reflection positivity ensures that the continuous-time spatial transfer operator is self-adjoint, enabling spatial scale doubling dual to imaginary time. We formalize the method as a general criterion for one-dimensional translation-invariant lattice systems and classify standard spin Hamiltonians according to their reflection positivity.
\end{abstract}
\maketitle
\setlength{\abovedisplayskip}{5pt plus 2pt minus 1pt}
\setlength{\belowdisplayskip}{5pt plus 2pt minus 1pt}

\appendix
\setcounter{secnumdepth}{2}
\setcounter{section}{3}

\section{Prerequisites of the Two-Purity Iteration}\label{app:uses}
The scale-doubling iteration in Ref.~\cite{main} comprises three elementary operations:
\begin{itemize}
\item \emph{Coupling identities.} Equalities between partition-function ratios on the spacetime torus, such as $\Zt42/\Ztsq[4]21$. These hold universally for every translation-invariant lattice Hamiltonian.
\item \emph{Temporal squaring.} The relation $\tncolumn{\rho_n,\rho_n}\propto\rho_n(2\beta)$ requires only positivity of the Gibbs operator, $e^{-\beta H}=e^{-\beta H/2}(e^{-\beta H/2})^\dagger\ge0$, which holds generally by Euclidean time reflection.
\item \emph{Spatial squaring.} The dual relation $\tnchain{\sigma_n,\sigma_n}\propto\sigma_{2n}$ requires positivity of the spatial transfer operator, $\sigma_n=X_\beta^n/Z_n\ge0$. This holds whenever $X_\beta$ is self-adjoint ($X_\beta=X_\beta^\dagger$), which additionally guarantees that $Z_n=\sum_i\lambda_i^n$ with real eigenvalues $\lambda_i$, permitting $\ell^p$ norm interpolation to non-dyadic intermediate lengths.
\end{itemize}
Self-adjointness of $X_\beta$ represents the spatial counterpart of Euclidean time reflection, and is not automatic. Once it holds, the remaining doubling argument follows directly.
The three hypotheses are therefore:
\begin{itemize}
\item[(H1)] \emph{Locality.} A translation-invariant chain with finite local
dimension and nearest-neighbour terms. Finite-range interactions reduce to
this by blocking sites.
\item[(H2)] \emph{Bond reflection positivity} (Sec.~\ref{app:rp}).
\item[(H3)] \emph{Input.} Both purities are near one at a single scale
(Sec.~\ref{app:criterion}).
\end{itemize}
Hypotheses (H1) and (H2) are algebraic properties of the Hamiltonian, whereas (H3) reflects the underlying spectral gap.

\section{Reflection positivity}\label{app:rp}
Let $\Theta$ be an antiunitary map on one site with $\Theta^2=\pm1$, for
example complex conjugation $K$ in a chosen basis. For an operator $C$ on one site write
$\Theta(C)=\Theta C\Theta^{-1}$.
\begin{itemize}
\item \emph{Reflection.} Cut the chain between two neighbouring sites. Mirror
the chain across the cut and apply $\Theta$ on every site. Call the result
$\theta$. It is antilinear, it maps the operators on the left half to those
on the right half, and $\theta^2=1$.
\item \emph{Bond form.} A nearest-neighbour Hamiltonian is in \emph{bond
reflection-positive form}~\cite{dls,fils,biskup,jaffejanssens} if it reads
\begin{equation}
 H=\sum_jD_j-\sum_{j,k}c_k\,C_{k,j}\,\Theta(C_k)_{j+1},\qquad c_k>0,
 \label{eq:rpform}
\end{equation}
with $\Theta(D)=D$. Each bond event of weight $c_k$ inserts $C_k$ on one
site and its mirror image $\Theta(C_k)$ on the next site,
\begin{equation*}
 -h_{j,j+1}=\sum_kc_k\;\bondform{C_k}{\Theta(C_k)}.
\end{equation*}
The sign in front matters. A term with $c_k<0$ cannot be absorbed by
rescaling $C_k$, because $\Theta$ conjugates the factor $i$:
$(iC)\otimes\Theta(iC)=C\otimes\Theta(C)$.
\end{itemize}
The model of Ref.~\cite{main} has this form with $\Theta=K$ in the Cartesian
basis and $c_k=1$. Its generators $G_\alpha=iS^\alpha$ are real, so
$-\bm S\cdot\bm S=\sum_\alpha G_\alpha\otimes\overline{G_\alpha}$.

\emph{Reflection positivity and transfer matrices.} Expand $e^{-\beta H}$ in the bond
events that cross the mirror. Each term is a left half followed by its mirror
image, so for every operator $A$ on the left half,
\begin{equation}
 \Tr\bigl[A\,\theta(A)\,e^{-\beta H}\bigr]=\rpcut
 =\sum\langle\Phi_A|\Phi_A\rangle\ge0 .
 \label{eq:rp}
\end{equation}
Reflection positivity began as the Osterwalder--Schrader axiom that
rebuilds a Hilbert space and a positive transfer matrix from Euclidean
data~\cite{os}. On the lattice it gives a self-adjoint, positive transfer
matrix~\cite{luscher}, and its Cauchy--Schwarz structure drives infrared
bounds and chessboard estimates~\cite{fss,fils,dls}. For spins and
Majoranas, Jaffe and Janssens characterize exactly which couplings are
reflection positive~\cite{jaffejanssens}. The two-purity method needs only
the weaker consequence that follows: the network cut at a single bond is
Hermitian. This is the Lüscher statement with space and time exchanged.

\begin{lemma}[Mirror is conjugation]\label{lem:mirror}
Under Eq.~\eqref{eq:rpform}, build the continuous-time MPO as in
Ref.~\cite{main}. The events are $(\tau,k)$, they occur at rate $c_k$, and
$e^{-\tau D}$ propagates between them. Then the transfer kernel obeys
\begin{equation}
 \overline{\Xab ab}=\Xab ba,\qquad\text{i.e.}\quad
 \overline{X_\beta(a,b)}=X_\beta(b,a).
 \label{eq:mirror}
\end{equation}
\end{lemma}
\begin{proof}
Open the vertex in imaginary time. Events on the left bond $a$ insert
$\Theta(C)$, and events on the right bond $b$ insert $C$, along a single traced
physical loop:
\begin{equation*}
 \overline{X_\beta(a,b)}=\overline{\eventsite{\Theta C}{C}}
 =\eventsite{C}{\Theta C}=X_\beta(b,a).
\end{equation*}
Write $\Theta=VK$. Complex conjugation of a trace conjugates every
insertion and keeps the time order. Conjugating the loop by $V$ then maps
$\overline{\Theta(C)}\mapsto C$, $\overline C\mapsto\Theta(C)$ and
$\overline D\mapsto D$, because $V\overline V=\pm1$ and $\Theta(D)=D$. The
picture that results is the mirror image of the original: the left
insertions are now $C$ and the right ones are $\Theta(C)$.
\end{proof}

\begin{theorem}[Space behaves like time]\label{thm:space}
Assume (H1) and (H2). Then $X_\beta$ is a self-adjoint Hilbert--Schmidt
operator on $L^2(\mu_\beta)$. For every even $n$, $Z_n(\beta)=\sum_i\lambda_i(\beta)^n$ with
real $\lambda_i$, and $\sigma_n(\beta)\ge0$. In particular the squaring
lemma, the coupling identities and the $\ell^p$ interpolation of
Ref.~\cite{main} hold verbatim.
\end{theorem}
\begin{proof}
The event measure $\mu_\beta$ is positive because $c_k>0$, and it is finite,
with mass $e^{\beta\sum_kc_k}$ per bond. Normalize $\|C_k\|=1$. The kernel is
then bounded, $|X_\beta(a,b)|\le d\,e^{\beta\|D\|}$, hence Hilbert--Schmidt,
and it is Hermitian by Lemma~\ref{lem:mirror}. The remaining statements are
those of Ref.~\cite{main} for a self-adjoint $X_\beta$.
\end{proof}

\emph{Relation to the sign problem.} A Hamiltonian is stoquastic if, in some basis, every
configuration of a quantum Monte Carlo expansion has a nonnegative weight.
Reflection positivity asks less of individual configurations and more of their
structure: a weight may be negative, but each configuration paired with
its mirror image contributes a square $|\cdot|^2$.
\begin{itemize}
\item In Ref.~\cite{main} the event measure is positive, but the entries of
$W$ and $X_\beta$ carry signs, because the $G_\alpha$ are real antisymmetric
matrices. A sampler over histories would meet a sign problem; the spectral
argument does not, since $\sum_i\lambda_i^n>0$.
\item Neither property implies the other. The quantum Heisenberg ferromagnet
is stoquastic but not reflection positive~\cite{speer}; see also
Proposition~\ref{prop:xyz}.
\item The sign-free fermion Monte Carlo methods based on time reversal or
Majorana symmetry exploit the same mirror-squared factorization,
$\det M\cdot\overline{\det M}$~\cite{wuzhang,wei,li,jaffe}.
\end{itemize}

\section{The general criterion}\label{app:criterion}
Define the one-step loss
\begin{equation}
 \delta(\varepsilon)=\tfrac12\bigl[(1-\varepsilon)^{-3}-1\bigr]^2 .
 \label{eq:delta}
\end{equation}
It lies below the diagonal, \deltacurve, up to $\varepsilon_*\simeq0.1273$,
the root of $\delta(\varepsilon)=\varepsilon$. Ref.~\cite{main} uses the
cruder bound $\delta\le5\varepsilon^2$, valid for $\varepsilon\le1/100$.

\begin{theorem}[Two-purity criterion]\label{thm:criterion}
Assume (H1) and (H2), and suppose that for some even $n_0$ and some $\beta_0>0$,
\begin{equation}
 \Ps(n_0,\beta_0),\ \Pt(2n_0,\beta_0)\ge1-\varepsilon_0,\qquad
 \varepsilon_0<\varepsilon_* .
 \label{eq:H3}
\end{equation}
Set $\varepsilon_{j+1}=\delta(\varepsilon_j)$, $n_j=2^jn_0$ and
$\beta_j=2^j\beta_0$. Then for every even $L\ge n_0$ the ground state
is unique and
\begin{equation}
 \Delta_L\ \ge\ c:=\inf_{j\ge0}\frac1{\beta_j}
 \log\frac{1-3\varepsilon_j}{3\varepsilon_j}\ >\ 0 .
 \label{eq:c}
\end{equation}
\end{theorem}
\begin{proof}
By Theorem~\ref{thm:space}, the four moves of the doubling estimate in
Ref.~\cite{main} apply without change. Move~1 gives
$\Ps(n,2\beta)\ge(1-\varepsilon)^3$. Move~2 squares in space to get
$\Ps(2n,2\beta)\ge1-\delta(\varepsilon)$. Move~3 gives
$\Pt(4n,\beta)\ge(1-\varepsilon)^2(1-\delta)\ge(1-\varepsilon)^3$, which
requires $\delta\le\varepsilon$, i.e.\ $\varepsilon\le\varepsilon_*$.
Move~4 squares in time to get $\Pt(4n,2\beta)\ge1-\delta(\varepsilon)$.
Induction on $j$ gives
$\Ps(n_j,\beta_j),\Pt(2n_j,\beta_j)\ge1-\varepsilon_j$, with
$\varepsilon_j$ decreasing because $\delta(\varepsilon)<\varepsilon$ on
$(0,\varepsilon_*)$.
Since $3\varepsilon_j<1/2$, the interpolation and gap-recovery steps of
Ref.~\cite{main} give uniqueness and $e^{-\beta_j\Delta_L}\le
3\varepsilon_j/(1-3\varepsilon_j)$ whenever $n_j\le L\le2n_j$.
Every term of the infimum is positive. Since
$\delta(\varepsilon)\le K\varepsilon^2$ near zero, $\varepsilon_j$ falls
doubly exponentially. The terms therefore converge to a positive limit, and
the infimum is positive.
\end{proof}
With $(n_0,\beta_0,\varepsilon_0)=(2304,784,1/100)$ this is the theorem of
Ref.~\cite{main}, with $c\ge(\log20)/784$.

\begin{corollary}[What (H3) excludes]\label{cor:obstruct}
Under (H1) and (H2), Eq.~\eqref{eq:H3} cannot hold at any scale
\begin{itemize}
\item[(i)] if the ground state is $k$-fold degenerate at the length $2n_0$
with $k\ge2$, because $\Pt\le w_0\le1/k$ there, a spectrum like \specbars{.49,.49,.02};
\item[(ii)] if $\liminf_L\Delta_L=0$ along even $L$, because
Theorem~\ref{thm:criterion} would contradict it.
\end{itemize}
\end{corollary}
Corollary~\ref{cor:obstruct} rules out gapless chains, symmetry-broken
chains, and chains whose ground state on the ring is asymptotically
degenerate.

The converse is expected but not proved. In a gapped chain with a unique
ground state and correlation length $\xi$, one expects
$1-\Ps\sim e^{-n/\xi}$ and $1-\Pt\sim n\,e^{-\beta\Delta}$. Both vanish
along the doubling sequence \torusseq, so Eq.~\eqref{eq:H3} should hold at
some finite scale, though perhaps a very large one.

For a conformal critical chain the picture is different. The extensive free
energy cancels in each torus ratio, so $\Ps$ and $\Pt$ approach ratios of
the universal torus partition function at a fixed aspect ratio. These are
numbers below one, so the defect stalls instead of collapsing.

\section{Classification of Reflection-Positive Hamiltonians}\label{app:catalogue}
\emph{Bilinear spin couplings.} Consider spin~$S$ with a nearest-neighbour
coupling matrix $J$ (real symmetric) and a uniform field $\bm h$,
\begin{equation}
 H=\sum_j\bm S_j^{\mathsf T}J\,\bm S_{j+1}+\bm h\cdot\bm S_j .
 \label{eq:bilinear}
\end{equation}
\begin{proposition}\label{prop:xyz}
Suppose there is a symmetric orthogonal $O$ with $\det O=-1$ such that
$-JO\succeq0$ and $O\bm h=\bm h$. Then Eq.~\eqref{eq:bilinear} is in the form
of Eq.~\eqref{eq:rpform}, for every $S$. For diagonal
$J=\operatorname{diag}(J_x,J_y,J_z)$ with nonzero entries and $\bm h=0$, such
an $O$ exists exactly when an even number of $J_x,J_y,J_z$ are negative
(ferromagnetic).
\end{proposition}
\begin{proof}
$O=-R$ with $R$ the identity or a $\pi$ rotation. Let $T$ be time reversal,
so $TST^{-1}=-S$, and put $\Theta=U_RT$. Then $\Theta\bm S\Theta^{-1}=O\bm S$.
Since $U_R$ commutes with $T$, $\Theta^2=U_R^2T^2=\pm1$. Factor
$-JO=\sum_kv_kv_k^{\mathsf T}$ with real $v_k$ and set
$C_k=v_k\cdot\bm S$. Then
$\sum_kC_k\otimes\Theta(C_k)=\sum\bm S^{\mathsf T}(-JO)O\bm S
=-\bm S^{\mathsf T}J\bm S$, and $\Theta(\bm h\cdot\bm S)=(O\bm h)\cdot\bm S$.
For diagonal $J$ the polar factor is unique: $O=\operatorname{sign}(-J)$,
and $\det O=-1$ counts the ferromagnetic entries modulo two.
\end{proof}
\noindent The proposition sorts the standard spin chains as follows.
\begin{itemize}
\item \emph{Reflection positive:}
 \begin{itemize}
 \item the Heisenberg antiferromagnet for every $S$, with $O=-1$;
 \item XXZ chains $S^xS^x+S^yS^y+\Delta S^zS^z$ with $\Delta\ge0$;
 \item the XY model with any signs, since a zero entry frees $O_{zz}$;
 \item the transverse-field Ising chain with either sign, with $O$ fixing
 the field axis.
 \end{itemize}
\item \emph{Not of this form:}
 \begin{itemize}
 \item the Heisenberg ferromagnet (three ferromagnetic entries);
 \item XXZ chains with $\Delta<0$;
 \item the antiferromagnet in a uniform field, since $O=-1$ reverses $\bm h$.
 \end{itemize}
\end{itemize}

The sublattice rotations by $\pi$ that map antiferromagnets to ferromagnets
change two signs at once, so they never alter this parity.

The uniform-field obstruction explains a choice made in Ref.~\cite{main},
Appendix~A. There the boundary chain is built with opposite endpoint fields
\fieldopp, since $\theta(-hS^z_1)=+hS^z_n$ makes that chain mirror
symmetric. The equal-field chain \fieldsame{} is reached only through the
symmetry $U_P$.

\emph{Spin-1 bilinear--biquadratic chain.}
Consider $H=\sum_j\cos\vartheta\,\bm S_j\cdot\bm S_{j+1}
+\sin\vartheta\,(\bm S_j\cdot\bm S_{j+1})^2$. In the Cartesian basis the
quadrupoles $Q_{\alpha\beta}=S^\alpha S^\beta$ are real. Their antisymmetric
part is $\frac12\epsilon_{\alpha\beta\gamma}G_\gamma$, so
\begin{equation}
 -h=\bigl(\cos\vartheta-\tfrac12\sin\vartheta\bigr)\sum_\alpha G_\alpha\otimes\overline{G_\alpha}
 -\sin\vartheta\sum_{\alpha,\beta}Q^{+}_{\alpha\beta}\otimes\overline{Q^{+}_{\alpha\beta}},
 \label{eq:blbq}
\end{equation}
where $Q^+_{\alpha\beta}=\frac12(Q_{\alpha\beta}+Q_{\beta\alpha})$; the cross terms cancel between symmetric and antisymmetric parts.
\begin{proposition}\label{prop:blbq}
With $\Theta=K$ in the Cartesian basis, the bilinear--biquadratic chain is
in reflection-positive form if and only if
$-\pi+\arctan2\le\vartheta\le0$.
\end{proposition}
\begin{proof}
The condition holds if both coefficients in Eq.~\eqref{eq:blbq} are
nonnegative, which is the stated window. For the converse, the form exists exactly
when the realigned bond matrix becomes positive semidefinite after some
on-site shift. Maximizing its least eigenvalue over all shifts gives a negative
value outside the window, for example $-0.087$ at $\vartheta=-0.66\pi$ and
$-0.156$ at $\vartheta=0.05\pi$ (Sec.~\ref{app:checks}).
\end{proof}

\begin{proofdiagram}
\begin{tikzpicture}[x=1cm,y=1cm,font=\footnotesize]
 \def\R{1.35}
 % reflection-positive window [-pi+arctan 2, 0], in degrees [-116.57, 0]
 \fill[mpored!14] (0,0)--(0:\R) arc[start angle=0,end angle=-116.57,radius=\R]--cycle;
 \draw[black!35,line width=.4pt] (0,0) circle[radius=\R];
 \draw[mpored,line width=2.2pt] (0:\R) arc[start angle=0,end angle=-116.57,radius=\R];
 % method applies: Haldane part of the window, (-45, 0]
 \draw[mpoblue,line width=2.2pt] (0:{\R+.12}) arc[start angle=0,end angle=-45,radius={\R+.12}];
 % phase boundaries
 \foreach \a in {45,90,-45,-135}{\draw[black!55,line width=.4pt] (\a:{\R-.12})--(\a:{\R+.12});}
 \node[black!70] at (-22:{\R+.75}) {Haldane};
 \node[black!70] at (67.5:{\R+.4}) {critical};
 \node[black!70] at (180:{\R+.45}) {FM};
 \node[black!70] at (-90:{\R+.38}) {dimerized};
 % special points
 \foreach \a in {0,18.43,45,-45,-116.57}{\fill (\a:\R) circle[radius=1.4pt];}
 \node[anchor=west,font=\scriptsize] at (18.43:{\R+.08}) {AKLT};
 \node[anchor=west,font=\scriptsize] at ([yshift=-1pt]0:{\R+.2}) {Heisenberg};
 \node[anchor=south west,font=\scriptsize] at (45:{\R+.02}) {ULS};
 \node[anchor=north west,font=\scriptsize] at (-45:{\R+.14}) {TB};
 \node[anchor=north east,font=\scriptsize] at (-116.57:{\R+.02}) {$-\pi{+}\arctan2$};
 \draw[flow,black!40] (0,0)--(35:.6) node[pos=1.15,black!60,font=\scriptsize] {$\vartheta$};
\end{tikzpicture}
\end{proofdiagram}
The red arc in the figure is the reflection-positive window. The blue arc,
$-\pi/4<\vartheta\le0$, is where it meets the Haldane phase, so
Theorem~\ref{thm:criterion} can apply there. Elsewhere in the window
Corollary~\ref{cor:obstruct} excludes it: the dimerized ground states are
degenerate, and the Takhtajan--Babujian point TB is critical.

The Heisenberg chain sits exactly on the edge of the window. The AKLT and
ULS points lie outside it: the positive biquadratic coefficient has the
wrong sign. Proving the AKLT gap this way would need a different $\Theta$, a
different cut, or blocking: a numerical search over every on-site
$\Theta$ finds none (Sec.~\ref{app:checks}).

\emph{Beyond the chain.} Finite-range couplings fit after blocking: put $r$
sites in a unit cell and reflect within the block. Reflection-positive
fermion models~\cite{jaffe} and two-dimensional lattices with reflection
planes~\cite{fils,biskup} supply the structural property (H2). In
two dimensions the spatial transfer operator acts on a whole column, so
(H3) becomes much harder to verify.

\section{Numerical checks}\label{app:checks}
The checks are short floating-point scripts kept outside this repository. Two tests are used.
\begin{itemize}
\item \emph{Gram test} (\texttt{rp\_check.py}). This tests
Eq.~\eqref{eq:rp} on open chains of $4$ and $6$ sites,
$\beta\in[0.3,3]$. The left-half operators $A$ index a Hermitian form, and
the test reports its least eigenvalue relative to the largest, minimized over
the on-site antiunitary $\Theta=VK$.
\item \emph{Realignment test} (\texttt{t4.py}). This checks the form of
Eq.~\eqref{eq:rpform} directly, with the on-site shift optimized.
\end{itemize}
The results:
\begin{itemize}
\item \emph{Spin-$\tfrac12$ XYZ chain.} All eight sign patterns of
$(J_x,J_y,J_z)=(\pm1,\pm.7,\pm.4)$ follow the parity rule: those with zero
or two ferromagnetic entries give $\ge-10^{-12}$ after optimizing $V$, and
the other four give $-0.048$. The antiferromagnet in a uniform field gives
$-0.095$, and the transverse-field Ising chain gives $\ge-10^{-12}$.
\item \emph{Spin-1, Cartesian $\Theta=K$.} The Heisenberg antiferromagnet
and $-(\bm S\cdot\bm S)^2$ pass. The ferromagnet fails at $-0.39$ and AKLT
fails at $-0.04$. Over $\vartheta$, the test is nonnegative at the sampled points inside the
window of Proposition~\ref{prop:blbq}, $\vartheta/\pi\in[-.62,0]$, and negative at the sampled points
outside it, $\vartheta/\pi\in\{-1,-.88,-.75\}\cup[.12,1]$.
\item \emph{Spin-1, every on-site $\Theta=VK$.} Minimizing over $V\in U(3)$
with twelve restarts does not help: AKLT stays at $-0.040$, ULS at $-0.18$,
$\vartheta=-0.75\pi$ at $-0.15$, and the ferromagnet improves only to $-0.15$.
\end{itemize}
These tests are necessary conditions for a particular $\theta$. A failure
excludes that $\Theta$, not every possible proof.
\begin{thebibliography}{13}
\bibitem{main}
Companion note, \emph{A diagrammatic proof of the Haldane gap via spacetime
tensor networks}, \texttt{outputs/thermal-mpo-gap-note.tex} (2026).
\bibitem{os}
K. Osterwalder and R. Schrader,
\href{https://doi.org/10.1007/BF01645738}{Axioms for Euclidean Green's
functions}, Commun. Math. Phys. \textbf{31}, 83 (1973).
\bibitem{luscher}
M. L\"uscher, Construction of a selfadjoint, strictly positive transfer
matrix for Euclidean lattice gauge theories, Commun. Math. Phys.
\textbf{54}, 283 (1977).
\bibitem{fss}
J. Fr\"ohlich, B. Simon, and T. Spencer,
\href{https://doi.org/10.1007/BF01608557}{Infrared bounds, phase
transitions and continuous symmetry breaking}, Commun. Math. Phys.
\textbf{50}, 79 (1976).
\bibitem{jaffejanssens}
A. Jaffe and B. Janssens,
\href{https://doi.org/10.1007/s00220-015-2545-z}{Characterization of
reflection positivity: Majoranas and spins}, Commun. Math. Phys.
\textbf{346}, 1021 (2016).
\bibitem{dls}
F. J. Dyson, E. H. Lieb, and B. Simon,
\href{https://doi.org/10.1007/BF01106729}{Phase transitions in quantum spin
systems with isotropic and nonisotropic interactions},
J. Stat. Phys. \textbf{18}, 335 (1978).
\bibitem{fils}
J. Fr\"ohlich, R. Israel, E. H. Lieb, and B. Simon,
\href{https://doi.org/10.1007/BF01940327}{Phase transitions and reflection
positivity. I.}, Commun. Math. Phys. \textbf{62}, 1 (1978).
\bibitem{biskup}
M. Biskup, \emph{Reflection positivity and phase transitions in lattice spin
models}, in \emph{Methods of Contemporary Mathematical Statistical
Physics}, Lecture Notes in Math. \textbf{1970} (Springer, 2009).
\bibitem{speer}
E. R. Speer, Failure of reflection positivity in the quantum Heisenberg
ferromagnet, Lett. Math. Phys. \textbf{10}, 41 (1985).
\bibitem{jaffe}
A. Jaffe and F. L. Pedrocchi, Reflection positivity for Majoranas,
Ann. Henri Poincar\'e \textbf{16}, 189 (2015).
\bibitem{wuzhang}
C. Wu and S.-C. Zhang, Sufficient condition for absence of the sign problem
in the fermionic quantum Monte Carlo algorithm, Phys. Rev. B \textbf{71},
155115 (2005).
\bibitem{wei}
Z.-C. Wei, C. Wu, Y. Li, S. Zhang, and T. Xiang, Majorana positivity and the
fermion sign problem of quantum Monte Carlo simulations,
Phys. Rev. Lett. \textbf{116}, 250601 (2016).
\bibitem{li}
Z.-X. Li, Y.-F. Jiang, and H. Yao, Majorana-time-reversal symmetries: a
fundamental principle for sign-problem-free quantum Monte Carlo
simulations, Phys. Rev. Lett. \textbf{117}, 267002 (2016).
\end{thebibliography}
\end{document}
